% Transcription — fonds Alexandre Grothendieck, shelfmark 90, batch 2 (pages 21-24).
% Established from the facsimile archives/batches/90/batch-02.pdf (a copy of
% 90.pdf fetched through the project's relay, https://grothendieck-relay.onrender.com/source/90.pdf;
% PDF page k = archivists' page 20+k, no cover sheet), per the skill
% transcribe-grothendieck. The folder is twenty-four pages; this is the second
% and last batch (pages 1-20 are batch 1, transcribed separately and possibly
% in parallel, so the notation here was fixed from these four pages alone).
% Pass: Opus 5.5 (claude-opus-5-5), 2026-09-26 — first pass, unchecked against the pages by a human.
%
% Pages skipped: none.
%
% Pages summarised: 22 and 24, two leaves of the typescript of EGA IV, §18.6
% (henselisation), paginated by hand « IV-951 » and « IV-950 » and scanned
% upside down (they were read turned). Page 24 (IV-950) carries the end of
% the proof of 18.6.6 (parts (i), (ii), (iv), (vi)) and 18.6.7 (the
% henselisation of a semi-local ring as the product of the henselisations of
% its localisations at the maximal ideals); page 22 (IV-951) continues 18.6.7,
% then Prop. 18.6.8 (B tensor hA for B finite semi-local over A) and the
% beginning of Th. 18.6.9. The typed text is published and is not re-set:
% each page gets one French \note, and the ink interventions (substitutions,
% insertions, completed references, cancelled passages, a marginal critique on
% p. 22 using the « ½ » = semi- shorthand of hand.md) are listed in full. The
% interventions are in a thin black pen, not the blue felt pen of pp. 21 and
% 23; the « ½ » shorthand points to his hand, but the attribution is not
% established and the notes say so.
%
% Pages transcribed from his hand: 21 and 23, blue felt pen, fast drafting.
% They are the rectos of the two typed leaves and read as one continuous
% argument (21 ends « Considérons alors », 23 opens on « Ñ, c'est ... »).
% The formulas carry; a fair share of the connecting prose does not and is
% left \ill{}.
%
% His pagination: none on pp. 21 and 23. The typescript's own numbers
% (IV-950, IV-951) are recorded in the notes.
%
% What the batch is about. Despite the folder title « KG (X) » and the
% inventory's filing under « Géométrie et topologie combinatoire », these
% pages hold the end of a proof on G-equivariant quasi-coherent sheaves via
% graded modules: the functors M -> M~ and E -> Γ·(E) (underlined Γ, the
% graded module of sections) are adjoint and the composite E -> (Γ·(E))~ is
% isomorphic to the identity, so Γ· is fully faithful; a subsheaf F of E
% corresponds to a sub-A·-module Γ·(F) stable under G; the smallest
% G-stable sub-A·-module N of Γ·(E) containing Γ·(F) is graded (argument
% degree by degree through the M^i); N~ is an O_X-G-module with
% Γ·(F)~ = F ⊂ N~ ⊂ Γ·(F')~ = F', hence is the sheaf sought (the smallest
% G-stable F' between F and E); its formation commutes with flat base change
% S' -> S; and if F is of finite type, so is F' (reduce to S affine, pick n
% with F(n) generated by sections φ_i in Γ(S, Γ^n(F)), take the G-invariant
% submodule N^n they generate). He notes in the margin that the beginning of
% the proof must be written better. That is: the G-stable saturation of a
% subsheaf, the lemma behind the existence of equivariant coherent
% extensions — the kind of statement equivariant K-theory needs, though
% nothing on these pages names K-groups. The two typed versos are unrelated
% to it: EGA IV §18.6, on henselisation of semi-local rings.
%
% Apparatus count: 30 \ill{}, 10 \uncertain{}.
\documentclass[11pt,a4paper]{article}
% Le préambule commun à toutes les transcriptions du fonds Grothendieck.
%
% Il tient en une idée : **l'incertitude se compose, elle ne se résout pas.**
% Une transcription qui lisse un mot douteux a détruit la seule chose pour
% laquelle on la faisait — un lecteur doit pouvoir savoir, à chaque endroit,
% ce qui était sur le feuillet et ce qui a été ajouté. Les six macros qui
% suivent sont donc l'appareil critique, et non de la décoration.
%
% Les mêmes commandes sont reconnues par `scripts/render.mjs`, qui produit la
% vue de lecture du site. Une macro ajoutée ici doit l'être aussi là-bas,
% faute de quoi le rendu échouera bruyamment — ce qui est voulu : un rendu
% silencieusement faux est pire qu'un rendu absent.

\ProvidesPackage{grothendieck}[2026/08/08 Transcriptions du fonds Grothendieck]

\usepackage{amsmath,amssymb,amsthm}
\usepackage{mathrsfs}
\usepackage{stmaryrd}
% Les diagrammes commutatifs sont partout dans ce fonds : c'est la langue
% même de la géométrie algébrique de Grothendieck, et une transcription qui
% les rendrait en prose ne transcrirait rien.
\usepackage{tikz-cd}
% Français par défaut — c'est la langue des feuillets — mais l'anglais est
% chargé aussi : la traduction et le résumé sont des documents anglais, et la
% typographie française (espaces avant les deux-points) y serait une faute.
% Ces documents basculent par \selectlanguage{english} en tête de corps.
\usepackage[english,french]{babel}
\usepackage{fontspec}
\usepackage{geometry}
\geometry{a4paper,margin=25mm,marginparwidth=28mm}
\usepackage{marginnote}
\usepackage{xcolor}
% ulem, not soul. `soul`'s \st and \ul are built on a character-by-character
% scan that XeLaTeX's font machinery breaks: the document fails with an
% undefined control sequence rather than degrading. ulem does the same two jobs
% and is Unicode-engine safe. `normalem` stops it hijacking \emph, which the
% transcriptions use for Grothendieck's own emphasis.
\usepackage[normalem]{ulem}
\usepackage{hyperref}
\hypersetup{colorlinks=true,linkcolor=black,urlcolor=black,pdfborder={0 0 0}}

% --- Métadonnées du lot -----------------------------------------------------
% Reprises telles quelles par le script de rendu, qui les affiche en tête de
% la vue de lecture. Elles fixent ce que le fichier prétend couvrir : sans
% elles, une transcription flotte sans point d'attache dans un fonds de seize
% mille feuillets.

\newcommand{\folder}[1]{\gdef\grothFolder{#1}}
\newcommand{\batch}[1]{\gdef\grothBatch{#1}}
\newcommand{\leaves}[2]{\gdef\grothFirst{#1}\gdef\grothLast{#2}}
\newcommand{\foldertitle}[1]{\gdef\grothFoldertitle{#1}}
\gdef\grothFolder{?}\gdef\grothBatch{?}\gdef\grothFirst{?}\gdef\grothLast{?}\gdef\grothFoldertitle{}

% --- Le résumé --------------------------------------------------------------
%
% Il ouvre la lecture modernisée, sous le titre « Résumé », et il est composé
% en petit. Ce n'est pas un ornement : c'est le seul passage du document qui
% ne soit pas des mathématiques, et un lecteur doit pouvoir le reconnaître
% avant de l'avoir lu — puis le sauter s'il connaît déjà le sujet, ou s'y
% arrêter s'il ne le connaît pas. Le filet qui le suit marque la reprise du
% corps.

\newenvironment{resume}
  {\small\setlength{\parskip}{0.45em}}
  {\par\medskip\hrule\medskip}

% --- Mots-clés ---------------------------------------------------------------
%
% En anglais, en clôture du résumé de la lecture modernisée : le vocabulaire
% moderne sous lequel chercher ce que les feuillets construisent. Le manifeste
% du site (`npm run manifest`) les extrait pour étiqueter la cote entière —
% cette ligne est la source unique des tags, il n'y a pas de fichier à côté.
\newcommand{\keywords}[1]{\par\smallskip\noindent
  {\footnotesize\textsc{Keywords} --- \textit{#1}}\par}

% --- L'appareil critique ----------------------------------------------------

% Le numéro de feuillet, en marge. C'est l'ancre qui permet de régler un
% désaccord sur une formule en regardant *un* feuillet plutôt qu'une chemise
% de six cents. Le site s'en sert aussi pour tourner les pages du fac-similé
% quand on fait défiler la transcription.
\newcommand{\leaf}[1]{%
  \marginnote{\footnotesize\color{gray!70}\textsf{#1}}%
  \label{leaf:#1}\ignorespaces}

% Illisible : on ne devine pas. Un mot inventé qui se lit comme les autres est
% le pire résultat possible.
\newcommand{\ill}{\textcolor{red!60!black}{[\,\ldots\,]}}

% Lecture douteuse : le mot est donné, et signalé comme incertain.
\newcommand{\uncertain}[1]{\uline{#1}}

% Ajout de l'éditeur — une lettre manquante, un mot sous-entendu.
\newcommand{\add}[1]{\textcolor{blue!50!black}{[#1]}}

% Biffé par Grothendieck lui-même. Ce qu'il a rayé fait partie du document :
% une rature dit souvent où la pensée a changé de direction.
\newcommand{\struck}[1]{\sout{#1}}

% Note du transcripteur, sur l'état matériel du feuillet.
\newcommand{\note}[1]{\par\smallskip{\small\color{gray!60!black}\textit{[#1]}}\par\smallskip}

% Note marginale de Grothendieck, distincte du corps du texte.
\newcommand{\marginal}[1]{\par\smallskip\noindent\hspace{1em}%
  {\small\color{gray!60!black}#1}\par\smallskip}

% --- Titre ------------------------------------------------------------------

\AtBeginDocument{%
  \begin{center}
    {\large\bfseries Fonds Alexandre Grothendieck --- cote n\textsuperscript{o}~\grothFolder}\\[2pt]
    {\itshape\grothFoldertitle}\\[4pt]
    {\small feuillets \grothFirst--\grothLast\ (lot \grothBatch)}\\[2pt]
    {\footnotesize\color{gray!60!black}Transcription établie d'après le fac-similé numérisé
     par l'Université de Montpellier.\\
     Les lectures incertaines sont soulignées, les illisibles notées [\,\ldots\,],
     les ajouts entre crochets.}
  \end{center}
  \vspace{1em}}

\endinput


\folder{90}
\batch{2}
\pages{21}{24}
\dating{[à partir de 1964]}
\watermark{Édition de démonstration}
\foldertitle{KG (X) : notes manuscrites (s.d.)}

\begin{document}

\page{21}
$M \mapsto \widetilde{M}$ en sens inverse (ils sont adjoints \ldots) et le
composé \struck{est} $E \mapsto \widetilde{\underline{\Gamma}^{\cdot}(E)}$ est
\uncertain{isom.} à l'identité, ce qui montre que $\underline{\Gamma}^{\cdot}$
est \emph{pleinement fidèle}.
\note{« (ils sont adjoints \ldots) » est ajouté en interligne et rattaché par
un trait après « en sens inverse ».}

$F$ correspond à un sous-\struck{\ill{}}$A^{\cdot}$-module $\Gamma^{\cdot}(F)$
de $\Gamma^{\cdot}(E)$, \ill{} \uncertain{à priori} stable par action de $G$.
Si \struck{$E' \supset F$ tel que} $F'$ tel que $F \subset F' \subset E$,
$F'$ stable, \ill{} stable par $G$ dans $\Gamma^{\cdot}(F')$
\struck{$\Gamma^{\cdot}(F') \subset \Gamma^{\cdot}(E)$}, \uncertain{s'identifie}
\[
  \underline{\Gamma}^{\cdot}(F) \subset \underline{\Gamma}^{\cdot}(F')
  \subset \underline{\Gamma}^{\cdot}(E)
\]
est \ill{} stable par $G$. On en conclut qu'il y a un plus petit
sous-\ill{}-module $N$ de $\Gamma^{\cdot}(E)$, contenant
$\underline{\Gamma}^{\cdot}(F)$, et stable par $G$.
\note{un mot court est inséré au-dessus de « sous- », avec un trait de
renvoi ; il n'a pas été lu.}
On voit de suite que $N$ est un sous-module gradué
[\ill{} \ill{}. $N \cap \Gamma^{i}(E)$ contient $\underline{\Gamma}^{i}(F)$,
et est stable par $G$, donc contient le plus petit sous-module $M^{i}$ de
$\underline{\Gamma}^{i}(E)$ contenant $\underline{\Gamma}^{i}(F)$ et stable par
$G$, donc $N$ contient $\sum M^{i}$, d'où par minimalité
\uncertain{lui est ég.} ?]. De \ill{}, par fonctorialité de la construction de
$N$, on voit que $N$ est un sous-$A^{\cdot}$-module. Considérons
\uncertain{alors}

\page{22}
\note{feuillet dactylographié d'EGA IV, paginé à la main « IV-951 », au verso
des notes ; texte publié, non recomposé. Il porte la fin de (18.6.7) (le
hensélisé $^{h}\!A$ d'un anneau semi-local possède la propriété universelle de
(18.6.6, (ii)) pour les homomorphismes semi-locaux), la Proposition (18.6.8)
($B \otimes_A {}^{h}\!A$ est une $B$-algèbre semi-locale isomorphe à
$^{h}\!B$ pour $B$ semi-locale et entière sur $A$) avec sa démonstration, et
le début du Théorème (18.6.9). Les interventions à la plume fine, d'une main
que le « ½ » pour \emph{semi-} rapproche de la sienne sans que l'attribution
soit établie :
\begin{itemize}
  \item « homomorphisme local » et « anneau local » : \emph{local} est biffé
    deux fois et remplacé par \emph{semi-} en interligne ;
  \item « l'idéal maximal de $B$ » devient « l'idéal maximal $\mathfrak{n}$
    de $B$ » ;
  \item le passage « Enfin, $^{h}\!A$ possède aussi la propriété
    universelle \ldots\ en découle aussitôt » est encadré et barré de traits
    obliques, avec en marge : « vider, (la forme hybride ici ne mérite pas le
    nom de "propr. universelle") ou bien \ill{} \ill{} de "hom ½ locaux" dans
    des \ill{} \ill{} » ;
  \item « anneau local hensélien » reçoit la lettre $C$ ; « ayant les
    propriétés précédentes » devient « ayant la propriété précédente » ;
  \item en marge de « d'un problème universel », un mot relié par un trait,
    \uncertain{cavalier} ;
  \item le « t » final de « le soit » est complété à la main.
\end{itemize}}

\page{23}
\note{suite directe de la page 21, dont ce feuillet est le suivant ; la page 22, dactylographiée, s'intercale sans rapport avec l'argument.}
$\widetilde{N}$, c'est un $\underline{O}_X$-$G$-module
\struck{\ill{} $X$} \struck{contenu}
$\subset \widetilde{\underline{\Gamma}^{\cdot}(F')} = F'$, et contenant
$\widetilde{\underline{\Gamma}^{\cdot}(F)} = F$.
\note{« $\underline{\Gamma}^{\cdot}(F') =$ » est ajouté en interligne entre
« $\subset$ » et « $F'$ », ce dernier surchargé.}
C'est donc le \struck{$F'$ \ill{}} faisceau cherché. Par construction, sa
formation commute à toute extension plate $S' \to S$ de la base.
Enfin, si $F$ est de type fini, \struck{alors $\exists$} prouvons que le
module $F'$ l'est. On peut supposer $S$ affine. Mais $\exists$ \uncertain{$n$}
tel que \uncertain{$F(n)$} \struck{\ill{} \ill{}} engendré par ses sections,
\struck{soient des} donc par des sections
$\varphi_i$ ($1 \leq i \leq N$) $\in \Gamma(S, \underline{\Gamma}^{n}(F))$.
Considérons alors le sous-module
\struck{$\underline{N}^{n}$ de $\underline{\Gamma}^{\cdot}(E)$ engendré par les
$\varphi_i$. Alors $\widetilde{N}^{\cdot} \subset F$ est un \ill{} faisceau
\ill{} $F$, comme on \ill{} \ill{} de suite. On \ill{} \ill{} \ill{} répéter la
construction \uncertain{précédente} \ill{}}
\ill{}, \ill{} $G$-invariant de $\underline{\Gamma}^{n}(F)$ engendré par les
$\varphi_i$, soit $N^{n}$, et le sous-Module de $E$ qu'il définit.
\emph{On} \emph{vérifie} quand même que c'est $\underline{F}'$, ce qui montre
que $F'$ est de type fini\ldots\ldots
\note{le passage biffé est enfermé dans un cadre et barré de traits obliques ;
le cadre s'arrête avant « $G$-invariant ».}
\marginal{Rédiger mieux ce début de la démonstration, pour pouvoir y \ill{}
\ill{}\ldots}

\page{24}
\note{feuillet dactylographié d'EGA IV, paginé à la main « IV-950 », au verso
des notes et précédant le feuillet de la page 22 ; texte publié, non recomposé.
Il porte la fin de la démonstration de (18.6.6) — parties (i), (ii), (iv) et
(vi) — et (18.6.7), définition du hensélisé d'un anneau semi-local comme
produit $\prod_i {}^{h}(A_{\mathfrak{m}_i})$. Mêmes plume fine et même
réserve sur l'attribution que pour la page 22. Les interventions :
\begin{itemize}
  \item « l'anneau local de $S'^{(\lambda)}$ » devient « l'anneau local
    $\mathcal{O}_\lambda$ de $S'^{(\lambda)}$ » ;
  \item « Il y a \ldots\ un $A_\lambda$-homomorphisme » : \emph{donc} est
    ajouté, et « $\mathcal{O}_\lambda \to {}^{h}\!A$ » écrit en interligne à
    la place d'un mot biffé ;
  \item en marge de « $\mathrm{Hom.loc}({}^{h}\!A, B) = \varinjlim
    \mathrm{Hom.loc}(A_\lambda, B)$ », le signe $\simeq$ ; les flèches sous
    les $\varinjlim$ sont tracées à la main ;
  \item la référence à Bourbaki, \emph{Alg. comm.}, chap.~III, §~2, est
    complétée en « n°~8, cor.~3 du th.~1 » ;
  \item le passage « et si $A' = B_{\mathfrak{n}}$ est une $A$-algèbre locale
    strictement essentiellement étale \ldots\ au point fermé de
    $\mathrm{Spec}(A)$ » est mis entre crochets et barré de croix ;
  \item « alors de (17.4.1) » est biffé et remplacé par « de la remarque qui
    suit 18.6.2 » ; « l'homomorphisme \ldots\ est bijectif » devient « les
    homomorphismes $A \to A'_\lambda$ sont bijectifs » ;
  \item « C'est une $A$-algèbre semi-locale » devient « C'est un $A$-module
    fidèlement plat et une $A$-algèbre semi-locale » ;
  \item la seconde référence à Bourbaki, chap.~III, §~2, est complétée en
    « n°~13, prop.~18 ».
\end{itemize}}

\end{document}
